\section{Lecture 20: Residual Networks and Skip Connections}

\subsection{Motivation}

In previous lectures, we studied:

\begin{itemize}
    \item gradient instability (Lecture 16),
    \item activation functions and depth (Lecture 17),
    \item normalization methods (Lecture 18),
    \item regularization (Lecture 19).
\end{itemize}

\medskip

\noindent
These techniques improve training, yet a fundamental problem remains.

\medskip

\noindent
\textbf{Observation.}  
As networks become deeper, performance may degrade.

\medskip

\noindent
\textbf{Key question.}  
Why does increasing depth sometimes hurt performance, and how can we fix it?

\subsection{The Degradation Problem}

Empirically, deeper networks can have:

\begin{itemize}
    \item higher training error,
    \item worse optimization behavior,
\end{itemize}

even though they are more expressive.

\medskip

\noindent
\textbf{Paradox.}
\begin{itemize}
    \item A deeper model should be able to represent a shallow model
    \item Yet optimization fails to find such solutions
\end{itemize}

\medskip

\noindent
\textbf{Insight.}  
The issue is not expressivity, but optimization difficulty.

\subsection{Residual Learning}

Residual networks introduce skip connections.

\medskip

\noindent
Instead of learning:
\[
H(x),
\]
we learn a residual function:
\[
F(x) = H(x) - x.
\]

\medskip

\noindent
Thus:
\[
H(x) = x + F(x).
\]

\medskip

\noindent
\textbf{Residual block.}
\[
y = x + F(x; \theta).
\]

\medskip

\noindent
\textbf{Interpretation.}

\begin{itemize}
    \item The network learns a correction to the identity
    \item If identity is optimal, set $F(x) = 0$
\end{itemize}

\medskip

\noindent
\textbf{Benefit.}
\begin{itemize}
    \item Easier optimization
    \item Improved gradient flow
\end{itemize}

\medskip

\noindent
\textbf{Insight.}  
Learning residuals is easier than learning full transformations.

\subsection{Gradient Flow in Residual Networks}

Consider:
\[
y = x + F(x).
\]

\medskip

\noindent
\textbf{Gradient.}
\[
\frac{\partial L}{\partial x}
= \frac{\partial L}{\partial y}
\left( I + \frac{\partial F}{\partial x} \right).
\]

\medskip

\noindent
\textbf{Key observation.}

\begin{itemize}
    \item The identity term ensures gradients can flow directly
    \item Avoids vanishing gradients
\end{itemize}

\medskip

\noindent
\textbf{Insight.}  
Skip connections create direct gradient pathways.

\subsection{Dynamical Systems Perspective}

Residual networks can be interpreted as discretizations of continuous systems.

\medskip

\noindent
\textbf{Form.}
\[
h_{k+1} = h_k + F(h_k).
\]

\medskip

\noindent
This resembles:
\[
\frac{dh}{dt} = F(h),
\]

which is a differential equation.

\medskip

\noindent
\textbf{Interpretation.}

\begin{itemize}
    \item Depth corresponds to time steps
    \item Network computes a flow in representation space
\end{itemize}

\medskip

\noindent
\textbf{Insight.}  
Deep networks can be viewed as discrete dynamical systems.

\subsection{Very Deep Architectures}

Residual connections enable training of very deep networks:

\begin{itemize}
    \item hundreds or even thousands of layers
\end{itemize}

\medskip

\noindent
\textbf{Key components.}

\begin{itemize}
    \item skip connections
    \item normalization (e.g., batch normalization)
    \item suitable activations (e.g., ReLU)
\end{itemize}

\medskip

\noindent
\textbf{Empirical observation.}

\begin{itemize}
    \item deeper residual networks often perform better
    \item training becomes more stable
\end{itemize}

\medskip

\noindent
\textbf{Connection to modern models.}

\begin{itemize}
    \item residual connections are used in:
    \begin{itemize}
        \item convolutional networks
        \item transformers
    \end{itemize}
\end{itemize}

\subsection{A Unifying Perspective}

Residual networks address a fundamental challenge:

\begin{itemize}
    \item enabling depth without sacrificing trainability.
\end{itemize}

\medskip

\noindent
They combine:

\begin{itemize}
    \item architectural design (skip connections),
    \item normalization,
    \item activation functions.
\end{itemize}

\medskip

\noindent
\textbf{Core idea.}  
Introduce identity mappings to stabilize both forward and backward signal propagation.

\medskip

\noindent
\textbf{Key insight.}  
Depth becomes useful only when it is trainable.

\subsection{Outlook}

In this lecture, we studied:

\begin{itemize}
    \item the degradation problem,
    \item residual learning,
    \item gradient flow in skip connections,
    \item dynamical systems interpretation.
\end{itemize}

\medskip

\noindent
Residual networks mark a turning point in deep learning, enabling very deep architectures.

\medskip

\noindent
In the next lectures, we study specialized architectures that exploit structure in data:

\begin{itemize}
    \item convolutional networks for spatial data,
    \item sequence models and attention for sequential data.
\end{itemize}

\medskip

\noindent
\textbf{Guiding principle.}  
Architectural design is central to making deep learning effective at scale.